\section{An intersection that determines an exponent}
\label{sec:intersection}

For $X\in F\setminus C$, the projective line through $[1]$ and $[X]$
consists of the nonzero elements of $C+CX$, modulo $C^\times$.
To recover such lines from degree-$p$
subfields, we need to pass from membership in $C(X)$ to membership in
$C+CX$, possibly after taking a power.  The following intersection provides such a
condition for two $p$-independent elements at once.  Its uniqueness assertion
will allow the exponent to be compared between different pairs.

The use of differential equations to control such intersections, including
the power correction in positive characteristic, appears in
\citet[Sections~3.4--3.5]{CadoretPirutka2021}.  Here the degree-$p$
subextensions allow a direct proof using commuting Euler derivations.

\begin{proposition}[Intersection and exponent]
\label{prop:intersection}
Let $F$ be a field of characteristic $p$ with $[F:F^p]<\infty$, and put
$C=F^p$.  Let $X,Y\in F^\times$ be $p$-independent.  Suppose
\[
 U\in C(X)^\times,\qquad V\in C(Y)^\times,
 \qquad [U]\notin\Fp[X],\qquad [V]\notin\Fp[Y].
\]
If
\begin{equation}
\label{eq:intersection-witness}
 X\,C(Y/X)^\times\ \cap\ U\,C(V/U)^\times\ne\varnothing,
\end{equation}
then there is a unique integer $n$, $1\le n<p$, such that
\begin{equation}
\label{eq:intersection-conclusion}
 U^n\in C+CX^n,\qquad V^n\in C+CY^n.
\end{equation}
In fact, either inclusion in \eqref{eq:intersection-conclusion} determines
$n$ uniquely.  In each inclusion, both coefficients are nonzero.
\end{proposition}

The shape of \eqref{eq:intersection-witness} reflects the inclusion
$CX+CY\subseteq X\,C(Y/X)$, and its counterpart for $U,V$.
Thus a common nonzero vector of the two spans yields the required intersection.

\begin{proof}
Choose $Z$ in the intersection \eqref{eq:intersection-witness}.  Extend
$X,Y$ to a $p$-basis and let $D_X,D_Y$ be the commuting Euler derivations
of Lemma~\ref{lem:p-basis}.  We first translate the two descriptions of $Z$
into differential equations.  Since $U$ and $V$ are not scalar multiples
of monomials in $X$ and $Y$, respectively, the eigenvector description in
that lemma gives
\[
 D_XU\ne0,\qquad D_YV\ne0.
\]
Define
\[
 B_X=\frac{U}{D_XU},\qquad B_Y=\frac{V}{D_YV},
 \qquad A_X=B_X-1,\qquad A_Y=B_Y-1.
\]
Here $A_X\in C(X)\setminus C$ and $A_Y\in C(Y)\setminus C$:
if, for example, $B_X\in C$, then $D_XU=B_X^{-1}U$, making $U$
an eigenvector and hence a scalar multiple of a monomial in $X$.
In particular, $A_X,A_Y,B_X,B_Y$ are all nonzero.

The derivations
\[
 \delta=D_X+D_Y,
 \qquad \delta'=B_XD_X+B_YD_Y
\]
satisfy $\delta X=X$, $\delta Y=Y$, $\delta'U=U$, and $\delta'V=V$.
Thus $\delta$ kills $Y/X$, whereas $\delta'$ kills $V/U$; both kill $C$.
The two expressions for $Z$ therefore give
\[
 \delta Z=Z=\delta'Z.
\]
Set
\[
 q_X=\frac{D_XZ}{Z},\qquad q_Y=\frac{D_YZ}{Z}.
\]
Dividing the preceding equations by $Z$ and subtracting yields
\begin{equation}
\label{eq:intersection-q}
 q_X+q_Y=1,\qquad A_Xq_X+A_Yq_Y=0.
\end{equation}
Both $q_X$ and $q_Y$ are nonzero: if one vanished, these two equations
would force the other $A$ to vanish.

Since $\delta$ commutes with $D_X$ and $D_Y$, the equation $\delta Z=Z$
implies that $q_X,q_Y$ are killed by $\delta$:
\[
 \delta q_X=\frac{D_X(\delta Z)}{Z}
              -\frac{D_XZ\,\delta Z}{Z^2}=0,
 \qquad \delta q_Y=0.
\]
By \eqref{eq:intersection-q}, $A_X/A_Y=-q_Y/q_X$, so
$\delta(A_X/A_Y)=0$.  Since $D_YA_X=D_XA_Y=0$, we obtain
\begin{equation}
\label{eq:intersection-eigenvalue}
 \frac{D_XA_X}{A_X}=\frac{D_YA_Y}{A_Y}.
\end{equation}
The left side lies in $C(X)$ and the right side in $C(Y)$.
Their intersection is $C$, by the $p$-monomial basis.  Consequently
$A_X$ and $A_Y$ are eigenvectors of their respective Euler derivations
with the same eigenvalue.  The distinct eigenvalues on
$1,X,\ldots,X^{p-1}$, and on $1,Y,\ldots,Y^{p-1}$, now give
\begin{equation}
\label{eq:intersection-monomials}
 A_X=c_XX^m,\qquad A_Y=c_YY^m
\end{equation}
for $c_X,c_Y\in C^\times$ and one integer $m$, $1\le m<p$.
The eigenvalue cannot be zero because $A_X,A_Y\notin C$.

It remains to integrate \eqref{eq:intersection-monomials} inside the
one-variable subfields.  Put $n=p-m$.  The identity
$1+c_XX^m=B_X\ne0$ ensures that all terms in
\[
 \frac{D_X(U^n)}{U^n}
 =\frac{n}{1+c_XX^m}
 =\frac{D_X(X^n+c_XX^p)}{X^n+c_XX^p}
\]
are defined.  Thus $D_X$ kills
$U^n/(X^n+c_XX^p)\in C(X)^\times$.  Its kernel on $C(X)$ is $C$,
so this quotient belongs to $C^\times$.  Since $X^p\in C^\times$,
we have $U^n\in C+CX^n$ with both coefficients nonzero.  Replacing
$X,U,c_X,D_X$ by $Y,V,c_Y,D_Y$ proves the second inclusion.

Finally, suppose $U^r=h+jX^r$ with $h,j\in C$ and $1\le r<p$.
Both coefficients must be nonzero: otherwise
$r[U]\in\Fp[X]$, contrary to $[U]\notin\Fp[X]$.
Taking logarithmic derivatives gives
\[
 \frac{D_XU}{U}=\frac{jX^r}{h+jX^r},
 \qquad A_X=\frac{h}{jX^r}.
\]
The last expression is a nonzero scalar multiple of $X^{p-r}$,
since $X^p\in C^\times$.  Comparison with
\eqref{eq:intersection-monomials} and linear independence of the
$p$ monomials forces $p-r=m$, hence $r=n$.
The argument for the inclusion involving $V$ is identical.
\end{proof}

All hypotheses and conclusions of Proposition~\ref{prop:intersection}
are unchanged when $X,Y,U,V$ are multiplied by elements of $C^\times$.
In particular, its exponent depends on projective points, rather than
on the representatives used in the calculation.
